% Printed 438.
\par\clearpage
\noindent\makebox[0pt][l]{\fontsize{13}{16}\selectfont 438}\hfill{\small\textbf{HEIGHT MEASUREMENTS.}}\hfill\null\par\medskip
\begin{center}
\setlength{\tabcolsep}{4pt}
\begin{tabular}{l|rrr|rrrr}
&\multicolumn{3}{c|}{Latitude}&\multicolumn{4}{c}{\shortstack{[E.] longitude\\from Göttingen}}\\
Hannover, Ägydius tower & \(52^\circ\)&\(22'\)&\(16^{\prime\prime}{,}4\)&\(-\)&\(0^\circ\)&\(12'\)&\(13^{\prime\prime}{,}8\)\\
\quad---\quad Neustädter tower &52&22&\(22{,}6\)&\(-\)&0&12&\(52{,}8\)\\
\quad---\quad Markt tower &52&22&\(24{,}8\)&\(-\)&0&12&\(28{,}4\)\\
\quad---\quad Kreuz tower &52&22&\(30{,}7\)&\(-\)&0&12&\(37{,}8\)\\
Neustadt am Rübenberge &52&30&\(21{,}8\)&\(-\)&0&28&\(53{,}7\)\\
Celle, south castle tower &52&37&\(31{,}4\)&\(+\)&0&8&\(4{,}9\)\\
\quad---\quad north castle tower &52&37&\(32{,}9\)&\(+\)&0&8&\(4{,}0\)\\
\quad---\quad tower of the town church &52&37&\(34{,}2\)&\(+\)&0&8&\(16{,}6\)\\
Lüneburg, Johannis tower &53&14&\(59{,}2\)&\(+\)&0&28&\(11{,}7\)\\
\quad---\quad Michaelisthurm &53&15&\(5{,}5\)&\(+\)&0&27&\(29{,}5\)\\
Hamburg, Michaelisthurm &53&33&\(1{,}8\)&\(+\)&0&2&\(3{,}0\)
\end{tabular}
\end{center}
\par\noindent\hspace*{\parindent}I cannot give the names of the three Hildesheim towers at the moment.
The first is what is called a roof turret, the second is a lantern, and the third is needle-shaped.
All these determinations take as their origin the site of the \textsc{Reichenbach}
meridian circle in the Göttingen observatory, and the calculations are based on the
dimensions of the terrestrial spheroid found by \textsc{Walbeck}. For the latitude of
the Brockenhaus, von \textsc{Zach}'s astronomical determination differs by \(10''\)
(his observing site was a little to the south of the tower). This skilful observer
attributes the discrepancy solely to the imperfection of the \textsc{Lenoir} circle
he used. I leave the matter open, but am rather inclined to think that the large
mass of the Harz mountains still lying south of the Brocken, whereas flat country
begins immediately to the north, has some part in it. It would certainly be
interesting if new astronomical observations were made in future with a more
perfect instrument on the Brocken or at its northern foot.
\par\medskip\noindent\hspace*{\parindent}At each of my triangulation points I measured the zenith distances of all the other
points connected to it with a 12-inch repeating circle, so that their relative
heights follow with great precision. Before publishing them, however, I should
first like to establish, by a more reliable

% Printed 439.
\par\clearpage
\noindent\null\hfill{\small\textbf{HEIGHT MEASUREMENTS.}}\hfill\makebox[0pt][r]{\fontsize{13}{16}\selectfont 439}\par\medskip
\noindent connection with sea level than has been made so far, the absolute heights above it.
I can, however, give you all my results so far for terrestrial refraction in the
following table. I would only remark that by refraction I mean the difference between
the directions of the ray of light at its two endpoints. What some astronomers call
refraction is really only half the refraction.
\begin{center}
\renewcommand{\arraystretch}{1.16}
\begin{tabular}{|l|r|r|r|}\hline
\multicolumn{1}{|c|}{Endpoints}&\multicolumn{1}{c|}{\shortstack{\rule{0pt}{1pt}\\Curvature of\\the terrestrial\\arc}}&\multicolumn{1}{c|}{\shortstack{Observed\\refraction}}&\multicolumn{1}{c|}{Ratio}\\\hline
Lichtenberg-Falkenberg &\(2767^{\prime\prime}{,}2\)&\(+345^{\prime\prime}{,}3\)&\(+0{,}1248\)\\
Deister-Falkenberg &\(2282{,}7\)&\(330{,}5\)&\(0{,}1448\)\\
Hohehagen-Brocken &\(2235{,}3\)&\(325{,}6\)&\(0{,}1457\)\\
Lichtenberg-Garssen &\(1961{,}7\)&\(252{,}0\)&\(0{,}1285\)\\
Deister-Garssen &\(1950{,}4\)&\(240{,}5\)&\(0{,}1233\)\\
Hils-Brocken &\(1779{,}1\)&\(292{,}6\)&\(0{,}1645\)\\
Lichtenberg-Deister &\(1597{,}2\)&\(252{,}6\)&\(0{,}1581\)\\
Hohehagen-Hils &\(1529{,}3\)&\(243{,}5\)&\(0{,}1592\)\\
Breithorn-Wilsede &\(1409{,}5\)&\(153{,}6\)&\(0{,}1090\)\\
Brocken-Lichtenberg &\(1372{,}5\)&\(185{,}2\)&\(0{,}1350\)\\
Hils-Deister &\(1316{,}8\)&\(175{,}3\)&\(0{,}1331\)\\
Hils-Lichtenberg &\(1291{,}1\)&\(187{,}1\)&\(0{,}1450\)\\
Hauselberg-Wilsede &\(1232{,}5\)&\(137{,}2\)&\(0{,}1113\)\\
Meridian mark-Hils &\(1188{,}9\)&\(187{,}0\)&\(0{,}1573\)\\
Falkenberg-Wilsede &\(1168{,}0\)&\(150{,}7\)&\(0{,}1290\)\\
Falkenberg-Wulfsode &\(1139{,}9\)&\(132{,}1\)&\(0{,}1159\)\\
Falkenberg-Scharnhorst &\(958{,}3\)&\(122{,}4\)&\(0{,}1277\)\\
Garssen-Falkenberg &\(910{,}0\)&\(78{,}3\)&\(0{,}0861\)\\
Falkenberg-Breithorn &\(884{,}0\)&\(109{,}2\)&\(0{,}1235\)\\
Hauselberg-Wulfsode &\(770{,}1\)&\(56{,}1\)&\(0{,}0728\)\\
Wulfsode-Wilsede &\(738{,}9\)&\(+84{,}6\)&\(+0{,}1146\)
\end{tabular}
\end{center}

% Printed 440.
\par\clearpage
\noindent\makebox[0pt][l]{\fontsize{13}{16}\selectfont 440}\hfill{\small\textbf{CORRESPONDENCE.}}\hfill\null\par\medskip
\begin{center}
\renewcommand{\arraystretch}{1.16}
\begin{tabular}{|l|r|r|r|}\hline
\multicolumn{1}{|c|}{Endpoints}&\multicolumn{1}{c|}{\shortstack{\rule{0pt}{1pt}\\Curvature of\\the terrestrial\\arc}}&\multicolumn{1}{c|}{\shortstack{Observed\\refraction}}&\multicolumn{1}{c|}{Ratio}\\\hline
Falkenberg-Hauselberg &\(691^{\prime\prime}{,}7\)&\(+87^{\prime\prime}{,}4\)&\(+0{,}1264\)\\
Meridian mark-Hohehagen &\(538{,}6\)&\(111{,}9\)&\(0{,}2078\)\\
Observatory-Hohehagen &\(447{,}0\)&\(54{,}3\)&\(0{,}1215\)\\
Scharnhorst-Breithorn &\(363{,}0\)&\(11{,}2\)&\(0{,}0309\)\\
Garssen-Scharnhorst &\(343{,}9\)&\(+20{,}6\)&\(+0{,}0599\)\\
Breithorn-Hauselberg &\(203{,}6\)&\(-23{,}2\)&\(-0{,}1141\)\\
Observatory-Meridian mark &\(162{,}3\)&\(+36{,}8\)&\(+0{,}2265\)
\end{tabular}
\end{center}
\par\noindent\hspace*{\parindent}The mean of all the determinations, taking the lengths of the lines into account,
would be \(0{,}1306\). One sees that the anomalies are much greater at small distances
than at large ones.
\par\noindent\hspace*{\parindent}On sunny summer days, in my experience, the refraction of light passing close
above the Earth's surface in flat country is usually negative in the morning and early
afternoon hours. The air is then always shimmering so strongly that very precise measurements
are impossible. The heliotrope light is then often a comet-shaped disc, sometimes as much as
a minute of arc in diameter. As the air gradually becomes more favourable in the afternoon, the disc
grows more and more concentrated and finally becomes a fine, intensely bright point like a
fixed star.
\begin{center}
[2.]\par\medskip
[The coefficient of refraction derived from the height measurements of 1824.]\par\bigskip
\textsc{Gauss} to \textsc{Schumacher}.\quad Göttingen, 28~November 1824.
\end{center}
\par\noindent\hspace*{\parindent}Only today am I able to communicate to you the result for terrestrial refraction
from my measurements this year. To obtain the amplitudes precisely, I first had to
determine the longitudes and latitudes of my triangulation points, and to do that I
first had to calculate all the triangles after adjusting the angles.
